\chapter{Introdução}
\label{cap:introducao}

O sequenciamento de nova geração (NGS) revolucionou o estudo de
comunidades microbianas ao tornar viável a análise de milhões de
sequências de DNA em uma única corrida. Projetos científicos que antes
exigiam anos de cultivo laboratorial podem hoje ser conduzidos em
semanas, desde que haja infraestrutura computacional adequada para
processar e interpretar o volume de dados gerado. Nesse contexto,
pesquisadores da Universidade Federal dos Vales do Jequitinhonha e
Mucuri (UFVJM) conduzem três investigações paralelas sobre comunidades
microbianas do solo: o projeto \textbf{INOVAHERB}, voltado à
caracterização do micobioma de plantas medicinais por meio de
sequenciamento ITS; o projeto \textbf{Pós-Fogo}, que monitora a
recuperação da microbiota bacteriana (16S rRNA) após eventos de
queimada; e o projeto \textbf{Biorremediação}, que correlaciona
comunidades fúngicas e bacterianas com variáveis edáficas em solos
contaminados por metais pesados.

Embora os três projetos compartilhem metodologia de sequenciamento e
análise estatística semelhantes, cada equipe utilizava ferramentas
distintas, scripts R isolados e fluxos de trabalho não padronizados.
A ausência de uma plataforma centralizada resultava em retrabalho,
dificuldade de reprodutibilidade e barreiras de acesso para membros da
equipe sem experiência em programação. Este trabalho surge como
resposta direta a esse problema: o desenvolvimento de uma plataforma
web acadêmica — denominada \textit{Rizoma} — capaz de integrar
o processamento de sequências FASTQ, a execução automática de análises
estatísticas em R e a visualização interativa dos resultados, tudo
sob uma interface acessível a pesquisadores de biologia.

\section{Motivação}

O principal gargalo identificado no fluxo de trabalho dos três
projetos não era a análise em si, mas a infraestrutura necessária para
executá-la de forma repetível. Análises como \textit{Analysis of
Compositions of Microbiomes with Bias Correction} (ANCOM-BC2)
\cite{lin2022analysis}, \textit{Multivariable Association with Linear
Models} (MaAsLin2) \cite{mallick2021multivariable} e redes de
co-ocorrência via SpiecEasi \cite{kurtz2015sparse} exigem
pré-processamento cuidadoso dos dados de abundância, parametrização
específica por tipo de marcador molecular e integração com metadados
ambientais das amostras. Realizá-las manualmente, sem rastreabilidade
de parâmetros ou histórico de execução, compromete a reprodutibilidade
científica exigida em publicações.

Adicionalmente, o volume de dados — arquivos FASTQ comprimidos na
ordem de gigabytes por projeto — inviabiliza o processamento em
estações de trabalho locais. A escolha de uma arquitetura de
microsserviços com orquestração Kubernetes permitiu consolidar os
recursos computacionais disponíveis no laboratório em um cluster de
cinco nós, distribuindo a carga de trabalho de forma eficiente.

\section{Objetivos}

\subsection{Objetivo Geral}

Projetar, implementar e validar uma plataforma web para análise
integrada de micobioma e metataxonômica, capaz de gerenciar o ciclo
completo desde o upload de arquivos FASTQ até a visualização dos
resultados estatísticos, com suporte aos três projetos científicos
ativos na UFVJM.

\subsection{Objetivos Específicos}

\begin{enumerate}
  \item Definir uma arquitetura de microsserviços adequada ao perfil
        de carga dos análises bioinformáticas, contemplando fila de
        jobs assíncrona, armazenamento de objetos e busca analítica.

  \item Implementar um R Worker capaz de executar as análises
        DESeq2, ANCOM-BC2, MaAsLin2, SpiecEasi, Random Forest, GSEA,
        FUNGuild e PICRUSt2 a partir de artefatos phyloseq armazenados
        em objeto-storage.

  \item Desenvolver uma API REST no padrão CQRS com FastAPI, expondo
        endpoints para gerenciamento de projetos, amostras, jobs e
        resultados.

  \item Construir um frontend Next.js 14 com visualizações interativas
        (PCoA, volcano plot, redes microbianas) conectadas em tempo
        real via WebSocket.

  \item Implantar a solução em cluster Kubernetes (k3s) com
        autenticação institucional Google OAuth.

  \item Gerar o painel unificado de seis PCoAs — figura central do
        TCC — ao término dos três pipelines.
\end{enumerate}

\section{Estrutura do Trabalho}

O restante deste texto está organizado da seguinte forma. O
Capítulo~\ref{cap:revisao} apresenta a revisão bibliográfica sobre
análise de microbiomas, ferramentas bioinformáticas e padrões
arquiteturais adotados. O Capítulo~\ref{cap:metodos} descreve a
metodologia de desenvolvimento, as decisões de arquitetura e os
detalhes de implementação de cada componente. O
Capítulo~\ref{cap:resultados} apresenta os resultados obtidos,
incluindo métricas de desempenho dos pipelines e as visualizações
geradas. Por fim, o Capítulo~\ref{cap:conclusao} reúne as conclusões,
limitações identificadas e trabalhos futuros.
